\section{A kernel on differences at prime arguments}\label{sec:kernel}

The tuple laws prescribe a distribution of admissible increments along each
edge of the distinguished cycle.  We now realize its low-denominator Fourier
multipliers by a signed kernel supported on polynomial values at prime
arguments.  This will turn avoidance of those values into cancellation for a
pair of Fourier lifts.  The kernel comparison adapts the integer-argument
method of \cite{IntersectivePower}; the unit conditions in its local weights
and their realization using prime arguments are the additional ingredients
here.

\subsection{Parameters and the pair estimate}

Fix the following parameters, all depending only on $h$ and the tuple
construction:
\begin{equation}\label{eq:parameters}
\begin{split}
 &\eta=\frac{1}{10000k},\qquad
 c_m=c_m(k,\eta/4)\quad\text{from Proposition~\ref{prop:prime-minor-arcs}},
 \qquad \nu=\frac{\min(\eta,c_m/k)}{10},\\
 &\beta=\frac{\nu}{100},\qquad \zeta=\frac{\nu}{100},\qquad
 \tau=\frac{\beta\gamma}{1000k(d+1)},\qquad
 \sigma=\frac{\gamma\tau}{4}.
\end{split}
\end{equation}
Here $\eta$ controls the major arcs, $\nu$ is a common saving from prime
sums, $\beta$ controls the Fourier cutoffs, and $\zeta$ bounds the
auxiliary parameter $l$.  The smaller exponents $\tau$ and $\sigma$ are
reserved for the descent in Section~\ref{sec:descent}.
The parameter $\eta$ is distinct from the exceptional masses $\eta_{l,p}$.
For a scale $N$ put
\begin{equation}\label{eq:kernel-cutoffs}
 Q=N^\beta,\qquad H=N^{\beta/2},\qquad
 \mathcal P_N=\{p\text{ prime}:P\le p\le Q\}.
\end{equation}

Fix $l\ge1$, write $F=h_l$, and choose a directed cycle edge
$v_j\to v_{j+1}$.  For each $p\ge P$, let $\pi_p$ be the distribution on
$\F_p$ of $z_{v_{j+1}}-z_{v_j}$ under $\mu_{l,p}^\circ$.  It is
supported on admissible residues by Proposition~\ref{prop:tuple-laws}.
For $v\in\Z$ define
\[
 \Gamma_p(v/p)=\sum_{y\in\F_p}\pi_p(y)e(vy/p).
\]
For a nonzero character, the conditional mixing estimate
\eqref{eq:tuple-mixing}, applied to the second endpoint and paired with the
conjugate character at the first, gives $|\Gamma_p(v/p)|\le p^{-\gamma}$.
If $\xi$ has squarefree denominator supported on primes at least $P$, write
$\xi=\sum_{p\mid q(\xi)}\xi_p$ for its prime-denominator decomposition and
put $\Gamma(\xi)=\prod_{p\mid q(\xi)}\Gamma_p(\xi_p)$.  Thus
\begin{equation}\label{eq:edge-multiplier-bound}
 \Gamma(0)=1,\qquad |\Gamma(\xi)|\le q(\xi)^{-\gamma}.
\end{equation}
The multiplier may depend on $l$ and the chosen edge; all estimates below
are uniform in both.

\begin{proposition}[Cancellation for an ordered pair]\label{prop:pair-cancellation}
Suppose $N$ is sufficiently large, $l\le N^\zeta$, and
$D=M_0\prod_{p\in B}p\le N^\beta$, where $B$ is a finite set of primes at
least $P$.  Let $J_1,J_2:[N]\to\C$ have modulus at most one and be
supported on residue classes $a'_1,a'_2\pmod D$.  Assume that
$r=a'_2-a'_1\pmod D$ is admissible at $l$ and that
\[
 J_1(x)J_2(y)\ne0\quad\Longrightarrow\quad
 x<y\quad\text{and}\quad y-x\notin\mathcal D_l.
\]
Set
\[
 \Xi=\{\xi\in\Q/\Z:q(\xi)\le H,
       \ q(\xi)\text{ squarefree},\ (q(\xi),D)=1\}.
\]
Then, for the multiplier of any directed cycle edge,
\begin{equation}\label{eq:pair-cancellation}
 D\left|\sum_{\xi\in\Xi}\Gamma(\xi)
       \widehat{J_1}_{[N]}(-\xi)\widehat{J_2}_{[N]}(\xi)\right|
 \ll H^{-\gamma/2}.
\end{equation}
\end{proposition}

Since $M_0$ contains every prime below $P$, the multiplier is defined on all
of $\Xi$.  To prove the proposition, consider the model kernel
\begin{equation}\label{eq:model-kernel}
 k_0(n)=\frac DN\,
   \mathbf1_{\{1\le n\le N\}}\mathbf1_{\{n\equiv r\ (D)\}}
   \sum_{\xi\in\Xi}\Gamma(\xi)e(-\xi n).
\end{equation}
We will construct a kernel $k_1$ supported on $\mathcal D_l$ such that
$\sup_\alpha|\sum_n(k_0(n)-k_1(n))e(\alpha n)|\ll H^{-\gamma/2}$.
The congruence restriction
in $k_0$ will be imposed on polynomial arguments by a local weight; the
factors $\Gamma_p$ will be imposed by further weights of mean one.

\subsection{Weights on polynomial arguments}

For a periodic function of the argument $m$, its \emph{unit mean modulo
$b$} is the uniform average over the classes $m\pmod b$ satisfying
$(T_l(m),b)=1$.  There are exactly
\begin{equation}\label{eq:unit-mean-count}
 n_b=b\prod_{\substack{p\mid b\\p\nmid l}}(1-1/p)
     =\frac{\varphi(lb)}{\varphi(l)}
\end{equation}
such classes.  Indeed, a prime dividing $l$ excludes no argument because
$(r_l,l)=1$, and any other prime excludes exactly one class modulo that
prime.  These means factor over prime powers.  They are also consistent
under replacing $b$ by a multiple: each allowed class has the same number
of allowed lifts.

At a prime $p\mid D$, choose an argument witnessing admissibility of
$r\pmod {p^{E_p}}$, and denote the valuation of its derivative by $s_p$.
Lemma~\ref{lem:local-lifting} supplies a class modulo $p^{E_p-s_p}$ on
which the unit condition holds and $F$ is constantly $r$ modulo
$p^{E_p}$.  Uniform argument lifts in this class give uniform higher
lifts of that value.  Combining the chosen classes gives one class modulo
\[
 d_r=\prod_{p\mid D}p^{E_p-s_p}\mid D.
\]
Let $\rho_D$ be $n_{d_r}$ times its indicator.  Then $\rho_D$ has unit mean
one and $\|\rho_D\|_\infty\le D$.

For every prime $p\nmid D$, necessarily $p\ge P$, choose, for each $y$
with $\pi_p(y)>0$, one allowed regular argument $x_y\pmod p$ satisfying
$F(x_y)=y$.  Define a function $\psi_p$ modulo $p$ by
\[
 \psi_p(x_y)=n_p\pi_p(y),\qquad
 \psi_p(x)=0\quad\text{at all other classes}.
\]
It has unit mean one, and its weighted $F$-value distribution is precisely
$\pi_p$.  Put
\[
 B_p(m)=\psi_p(m)-1,\qquad
 B_u(m)=\prod_{p\mid u}B_p(m),\qquad B_1=1,
\]
where $u$ is squarefree and $(u,D)=1$.  Each $B_p$ has unit mean zero,
$B_u$ has period $u$, and $|B_u|\le u$.  In particular these definitions
specify the weights also on nonunit arguments; those arguments will have
zero main term in the prime-progression estimate.

Expanding a finite product of $\psi_p=1+B_p$ produces a sum of the $B_u$
over squarefree products of those primes.  We retain $u\le H$ to bound
the periods of the weights.  In a unit average against a rational phase
$e(\alpha F(m))$, any prime factor of $u$ absent from $q(\alpha)$ makes
that term vanish.  Thus, when the part of $q(\alpha)$ supported away from
$D$ is at most $H$, the truncation retains every potentially nonzero term
of the corresponding product expansion.  The comparison below will also
bound the contribution from larger denominators.

By Lemma~\ref{lem:auxiliary}, all coefficients of $F$ are $O_h(a_l)$.
Consequently there is an integer $x_*\ge1$, independent of $l$, such that
for $x\ge x_*$,
\begin{equation}\label{eq:kernel-polynomial-scale}
 F(x)\asymp a_lx^k,\qquad F'(x)\asymp a_lx^{k-1},\qquad
 |F''(x)|\ll a_lx^{k-2},
\end{equation}
with $F,F'$ positive.  Put $U=(N/a_l)^{1/k}$.  Since $a_l$ is a positive
integer and $a_l\le al^k$,
\[
 U\le N^{1/k},\qquad U\gg N^{1/k-\zeta}.
\]
For large $N$ there is therefore a unique $X\ge x_*$ with $F(X)=N$, and
$X\asymp U$.  With $\Lambda'(n)=(\log n)\mathbf1_{\{n\text{ prime}\}}$,
define the signed kernel below.  The factor $\varphi(l)/l$ compensates
for the prime density in the progression $T_l(m)$, and $F'(m)/N$ turns
the resulting $m$-integral into normalized measure on the $F$-values:
\begin{equation}\label{eq:prime-kernel}
\begin{split}
 k_1(n)=\sum_{x_*<m\le X}&\mathbf1_{\{n=F(m)\}}
 \frac{F'(m)}{N}\frac{\varphi(l)}{l}\Lambda'(T_l(m))\rho_D(m)\\[-2pt]
 &\hspace{15mm}\times
 \sum_{\substack{u\le H\\u\text{ squarefree}\\(u,D)=1}}B_u(m).
\end{split}
\end{equation}
Here and below the argument $m$ is an integer in sums.  Every nonzero
term has $0<F(m)\le N$ and $T_l(m)$ prime, so $k_1$ is supported on
$\mathcal D_l\cap[N]$.  No positivity is required of $k_1$.

Choose a fixed constant $C$ so that $Y=ClU$ bounds $T_l(X)$.  The relation
$a_l=al^k/\lambda(l)$ gives
\begin{equation}\label{eq:prime-kernel-range}
 Y\gg N^{1/k},\qquad Y\ll N^{1/k+\zeta},\qquad
 Y^k\asymp N\lambda(l).
\end{equation}
In particular $Y\le N$ for sufficiently large $N$.

\subsection{Uniform Fourier comparison}

For $i=0,1$ use positive-sign transforms
$S_i(\alpha)=\sum_n k_i(n)e(\alpha n)$ on $\R/\Z$.
The signs differ from the negative-sign convention for Fourier
coefficients of functions; this choice makes $S_0$ concentrate at the
frequencies occurring in \eqref{eq:model-kernel}.

\begin{lemma}[Comparison of the two kernels]\label{lem:kernel-comparison}
For the data and ranges of Proposition~\ref{prop:pair-cancellation},
\begin{equation}\label{eq:kernel-comparison}
 \sup_{\alpha\in\R/\Z}|S_0(\alpha)-S_1(\alpha)|
 \ll H^{-\gamma/2}.
\end{equation}
\end{lemma}

\begin{proof}
Use major arcs of radius $3N^{-1+\eta}$ around the reduced rationals with
denominator at most $N^\eta$.  They are disjoint for large $N$, since
$3\eta<1$.  The remaining points form the minor arcs.  We first bound both
transforms on the minor arcs; on a major arc we will compare their local
coefficients.

Write $V_N(y)=N^{-1}\sum_{n\in[N]}e(yn)$.  Expanding the congruence
indicator in \eqref{eq:model-kernel} gives
\begin{equation}\label{eq:model-transform}
 S_0(\alpha)=\sum_{\xi\in\Xi}\sum_{j'\bmod D}
 \Gamma(\xi)e(-j'r/D)V_N(\alpha-\xi+j'/D).
\end{equation}
There are $O(DH^2)$ terms.  Their centers $\xi-j'/D$ are distinct:
an equality would give a rational with denominator both coprime to $D$
and dividing $D$, and hence an integer.  Each center has denominator at
most $DH\le N^{3\beta/2}\le N^\eta$.  On minor arcs, geometric summation
therefore gives
\begin{equation}\label{eq:model-minor}
 |S_0(\alpha)|\ll DH^2N^{-\eta}.
\end{equation}

To estimate $S_1$, first consider, with any linear twist $y'$, the sums
\[
 \sum_{m\in I}\Lambda'(T_l(m))e(\alpha F(m)+y'm)
\]
over subintervals $I$ of $(x_*,X]$.  On setting $n=T_l(m)$, expand the
restriction $n\equiv r_l\pmod l$ as the average of its $l$ additive
characters.  Each resulting prime sum has a polynomial phase in $n$ whose
coefficient of $n^k$ is $b_k=\alpha a/\lambda(l)$; the new twists change
only the linear coefficient because $k\ge2$.

Fix a representative of $\alpha$ and apply Dirichlet approximation to
$\vartheta=(k!)^2b_k$ with
\[
 Q_0=\lfloor Y^kN^{-\eta/2}\rfloor.
\]
The resulting reduced fraction $a'/q$ satisfies
$q\le Q_0$ and $|\vartheta-a'/q|\le1/(qQ_0)\le q^{-2}$.
By \eqref{eq:prime-kernel-range}, this places $\alpha$ within
$O(N^{-1+\eta/2})$ of a rational of denominator at most $(k!)^2aq$.
For a minor-arc point that denominator must exceed $N^\eta$ for large
$N$.  Hence, using $Y\le N$ and absorbing the fixed factor $(k!)^2a$,
\[
 Y^{\eta/4}\le q\le Y^{k-\eta/4}.
\]
Proposition~\ref{prop:prime-minor-arcs} applies uniformly in the linear
twists and in the interval, and gives a bound $O(Y^{1-c_m})$.
The relation $Y=ClU$ and the lower bound in
\eqref{eq:prime-kernel-range} imply
\[
 Y^{1-c_m}\ll U\,lN^{-c_m/k}
 \le U N^{\zeta-c_m/k}\ll U N^{-\nu}.
\]
This accounts for the dependence on $l$ before any periodic weight is
expanded.

By \eqref{eq:kernel-polynomial-scale}, the supremum and total variation of
$F'/N$ on $[x_*,X]$ are $O(U^{-1})$.  Partial summation therefore retains
an $O(N^{-\nu})$ bound after including this weight.  Use normalized
Fourier coefficients on each finite period for the other weights.  Since
$\rho_D$ is a constant times one residue-class indicator, its Fourier
$\ell^1$ norm is $n_{d_r}\le D$.  For $B_u$ the bound by its period times
its supremum gives Fourier $\ell^1$ norm at most $u^2$.  Expanding the
weights in \eqref{eq:prime-kernel}, and using $\varphi(l)/l\le1$, yields
\begin{equation}\label{eq:prime-minor-kernel}
 |S_1(\alpha)|\ll D\sum_{u\le H}u^2N^{-\nu}
 \ll DH^3N^{-\nu}.
\end{equation}
The powers in \eqref{eq:model-minor} and
\eqref{eq:prime-minor-kernel} are $2\beta-\eta$ and
$5\beta/2-\nu$, respectively.  Both are smaller than $-\beta\gamma/4$,
so these bounds are $O(H^{-\gamma/2})$.

Now let $\alpha=\alpha_0+y$ be on a major arc, with
$q=q(\alpha_0)\le N^\eta$ and $|y|\le3N^{-1+\eta}$.
Decompose $\alpha_0=\alpha_D+\alpha_1$ into the prime-power denominator
parts supported, respectively, on primes dividing $D$ and on primes not
dividing $D$.  The coefficient in \eqref{eq:model-transform} belonging
to the center $\alpha_0$, or zero if that center is absent, is
\begin{equation}\label{eq:model-local-coefficient}
 c_0(\alpha_0)=
 \begin{cases}
 e(\alpha_Dr)\Gamma(\alpha_1),
       &q(\alpha_D)\mid D\ \text{and}\ \alpha_1\in\Xi,\\
 0,&\text{otherwise}.
 \end{cases}
\end{equation}
Here the sign follows from $\alpha_D=-j'/D$ at that center.  Every other
center is at circle distance at least $1/(qDH)$ from $\alpha_0$.
Since $|y|qDH=o(1)$, geometric summation in each of the $O(DH^2)$ terms
gives
\begin{equation}\label{eq:model-major}
 S_0(\alpha)=c_0(\alpha_0)V_N(y)+O(qD^2H^3/N).
\end{equation}

We next derive the corresponding coefficient for $S_1$, including the
normalization of the prime weights.  For squarefree $u\le H$ coprime to
$D$, let $A_u$ be the unit mean of
\[
 b_u(m)=\rho_D(m)B_u(m)e(\alpha_0F(m)).
\]
Its period divides $b=qDu$, and $\|b_u\|_\infty\le Du$.
For $x_*\le x'<x''\le X$, Proposition~\ref{prop:prime-progressions}
gives
\begin{equation}\label{eq:weighted-prime-progression}
 \sum_{x'<m\le x''}b_u(m)\frac{\varphi(l)}l\Lambda'(T_l(m))
 =A_u(x''-x')+
 O\big(qD^2u^2Y^{31/32}\log^2(2lbY)\big).
\end{equation}
Indeed, one class $m\equiv c\pmod b$ corresponds to the class
$r_l+lc\pmod {lb}$ in a prime interval of length $l(x''-x')$.
It is reduced exactly when $(T_l(c),b)=1$, since $(r_l,l)=1$.
For such a class the factor $\varphi(l)/l$ changes the prime main term
to $(x''-x')/n_b$, by \eqref{eq:unit-mean-count}.
On a nonreduced class the only possible primes divide $lb$, and their
total logarithmic weight is $O(\log(lb))$, with zero main term.
Apply the progression bound at both endpoints, then sum its errors over
at most $b=qDu$ classes with weights of modulus at most $Du$.
This proves \eqref{eq:weighted-prime-progression}, including the
nonreduced classes.

The smooth factor $F'(x)e(yF(x))/N$ has supremum plus total variation
\[
 O\big((1+N|y|)/U\big).
\]
For example, its derivative is bounded in integral by the contributions
of $F''/N$ and $2\pi |y|(F')^2/N$, which are $O(U^{-1})$ and
$O(N|y|/U)$ by \eqref{eq:kernel-polynomial-scale}.
Partial summation in \eqref{eq:weighted-prime-progression} therefore
produces main term
\[
 \frac{A_u}{N}\int_{x_*}^{X}F'(x)e(yF(x))\,dx
 =\frac{A_u}{N}\int_{F(x_*)}^{N}e(yt')\,dt'.
\]
The normalized integral $N^{-1}\int_{F(x_*)}^{N}e(yt')\,dt'$ differs
from $V_N(y)$ by
$O(a_l/N+(1+N|y|)/N)$: the missing initial interval has length
$F(x_*)\ll a_l$, and comparison of the integral on $[0,N]$ with its
integer sum costs $O(1+N|y|)$.
It follows that
\begin{equation}\label{eq:prime-major}
 S_1(\alpha)=c(\alpha_0)V_N(y)+O(H^{-\gamma/2}),\qquad
 c(\alpha_0)=
 \sum_{\substack{u\le H\\u\text{ squarefree}\\(u,D)=1}}A_u.
\end{equation}
Here are the error bounds underlying this assertion.  Since
$Y^{31/32}/U\ll lN^{-1/(32k)}$, the summed progression error is at most
\[
 \frac{1+N|y|}{U}\,qD^2
 \sum_{u\le H}u^2Y^{31/32}N^{o(1)}
 \ll N^{2\eta+4\beta+\zeta-1/(32k)+o(1)}.
\]
The replacement of the integrals costs, using $|A_u|\le Du$,
\[
 O\big(DH^2(U^{-k}+N^{\eta-1})\big)
 \ll N^{2\beta+k\zeta-1}+N^{2\beta+\eta-1}.
\]
The error in \eqref{eq:model-major} is at most
$N^{\eta+7\beta/2-1}$.  All three estimates are
$O(N^{-\beta\gamma/4})$: indeed $\beta,\zeta\le\eta/1000$,
$\eta=1/(10000k)$, and $\gamma<1$.  The first exponent has a strict
margin, so its $o(1)$ loss is harmless.

It remains to compare $c$ and $c_0$.  This is precisely the reason for
the unit means and the regular argument classes used above.  At primes
dividing $D$, the common factor in each $A_u$ is
\begin{equation}\label{eq:kernel-D-factor}
 \chi_D(\alpha_0)=
 \begin{cases}
 e(\alpha_Dr),&q(\alpha_D)\mid D,\\
 0,&\text{otherwise}.
 \end{cases}
\end{equation}
To see both cases, factor $\rho_D$ and the unit mean prime by prime.
At $p\mid D$ the chosen class has weight of unit mean one and prescribed
$F$-value modulo $p^{E_p}$.  A character of conductor at most $p^{E_p}$
therefore gives its value at $r$.  A primitive character of higher
conductor averages to zero, since Lemma~\ref{lem:local-lifting} gives
uniform higher lifts of $F$ on that class, all still satisfying the unit
condition.

Put $q_1=q(\alpha_1)$.  If $u$ has a prime factor outside $q_1$, its
local $B_p$ factor has unit mean zero and $A_u=0$.
For $p^j\Vert q_1$, write $\alpha_p$ for the $p^j$-part of $\alpha_1$
and let
\[
 a_{p,j}=\mathbb E_{\substack{m\bmod p^j\\p\nmid T_l(m)}}
                  e(\alpha_pF(m)),\qquad
 b_{p,j}=\mathbb E_{\substack{m\bmod p^j\\p\nmid T_l(m)}}
                  B_p(m)e(\alpha_pF(m)).
\]
The surviving $A_u$, for squarefree divisors $u$ of $q_1$, are
$\chi_D(\alpha_0)$ times products that choose $b_{p,j}$ when $p\mid u$
and $a_{p,j}$ otherwise.  Moreover,
\begin{equation}\label{eq:kernel-local-realization}
 a_{p,j}+b_{p,j}=
 \begin{cases}
 \Gamma_p(\alpha_p),&j=1,\\
 0,&j\ge2.
 \end{cases}
\end{equation}
For $j=1$, the weight $1+B_p=\psi_p$ pushes the unit mean forward to
$\pi_p$.  For $j\ge2$, it is supported on allowed regular classes
modulo $p$, and their uniform argument lifts give uniform lifts of the
$F$-value.  The primitive character therefore has mean zero.

If $q_1\le H$, every squarefree divisor needed in the product expansion
occurs in the sum for $c$.  Thus \eqref{eq:kernel-local-realization}
and \eqref{eq:kernel-D-factor} show that $c=c_0$, including the case of
a nonsquarefree $q_1$, when both vanish.  If $q_1>H$, the model
coefficient is zero.  Lemma~\ref{lem:local-sums} gives
$|a_{p,j}|\le p^{-8\delta j}$, while
\eqref{eq:edge-multiplier-bound} and
\eqref{eq:kernel-local-realization} give
$|b_{p,j}|\le2p^{-\gamma j}$ (for $j\ge2$ one has $b_{p,j}=-a_{p,j}$).
Bounding the truncated product expansion by the full sum of absolute
values, and using \eqref{eq:prime-thresholds}, we obtain
\[
 |c(\alpha_0)|\le
 \prod_{p^j\Vert q_1}(|a_{p,j}|+|b_{p,j}|)
 \le\prod_{p^j\Vert q_1}3p^{-\gamma j}
 \le q_1^{-\gamma/2}<H^{-\gamma/2}.
\]
Indeed $p^{\gamma/2}\ge3$, so each factor is at most
$p^{-\gamma j/2}$.  Together with \eqref{eq:model-major} and
\eqref{eq:prime-major}, this proves the comparison on major arcs.
The minor-arc bounds already proved complete the lemma.
\end{proof}

\begin{proof}[Proof of Proposition~\ref{prop:pair-cancellation}]
Extend $J_1,J_2$ by zero to $\Z$ and put
\[
 \mathcal B_i=\frac1N\sum_{x,n\in\Z}
                    k_i(n)J_1(x)J_2(x+n),\qquad i=0,1.
\]
To keep track of signs and normalizations, let
$\widetilde J_i(\alpha)=\sum_xJ_i(x)e(-\alpha x)$, with counting
measure in this transform.  Character orthogonality gives
\[
 \mathcal B_i=\frac1N\int_{\R/\Z}
 S_i(\alpha)\widetilde J_1(-\alpha)
                 \widetilde J_2(\alpha)\,d\alpha.
\]
Both counting-measure $\ell^2$ norms are at most $\sqrt N$.
Cauchy--Schwarz, Plancherel, and Lemma~\ref{lem:kernel-comparison}
therefore imply
$|\mathcal B_1-\mathcal B_0|\ll H^{-\gamma/2}$.
The forbidden-support hypothesis gives $\mathcal B_1=0$.
For every contributing pair in $\mathcal B_0$, the ordering and the
residue classes imply $1\le n\le N$ and $n\equiv r\pmod D$.
Substituting \eqref{eq:model-kernel} and setting $y=x+n$ thus yields
\[
 \mathcal B_0=\frac{D}{N^2}\sum_{\xi\in\Xi}\Gamma(\xi)
       \sum_{x,y}J_1(x)J_2(y)e(-\xi(y-x))
 =D\sum_{\xi\in\Xi}\Gamma(\xi)
       \widehat{J_1}_{[N]}(-\xi)\widehat{J_2}_{[N]}(\xi).
\]
This is exactly \eqref{eq:pair-cancellation}.
\end{proof}
